\documentclass[../../main.tex]{subfiles}
\graphicspath{{../../image/}{../../figure/}}
\begin{document}

\section{熵--熵流对}\label{section:ch07-熵熵流对}

熵在一维双曲守恒律组的理论中扮演核心角色。把第 3.2 节的讨论适配到当前设定，我们推断 $\mathscr{O}$ 上的函数 $\eta$ 与 $q$ 构成方程组\eqref{align:ch07-典范守恒律}的一个\emph{熵--熵流对}，当
\begin{align}
\mathrm{D}q(U)=\mathrm{D}\eta(U)\mathrm{D}F(U),\qquad U\in\mathscr{O}.
\label{align:ch07-熵流对定义}
\end{align}
进而，可积性条件\eqref{align:ch03-可积性}
在此退化为
\begin{align}
\mathrm{D}^2\eta(U)\mathrm{D}F(U)=\mathrm{D}F(U)^\top\mathrm{D}^2\eta(U),\qquad U\in\mathscr{O}.
\label{align:ch07-可积性}
\end{align}

从左乘 $R_j(U)^\top$、从右乘 $R_k(U)$（$j\neq k$）于\eqref{align:ch07-可积性}，我们推出\eqref{align:ch07-可积性}等价于
\begin{align}
R_j(U)^\top\mathrm{D}^2\eta(U)R_k(U)=0,\qquad j,k=1,\cdots,n;\ j\neq k,
\label{align:ch07-熵可积正交}
\end{align}
其理解是：当 $\lambda_j(U)\neq\lambda_k(U)$ 时\eqref{align:ch07-熵可积正交}自动成立，但对关联于多重特征速度的特征向量 $R_i$ 可能需要重新规范化。（与\eqref{align:ch03-可积性投影}
相比较。）注意，某熵 $\eta$ 为凸的要求现在可方便地表达为
\begin{align}
R_j(U)^\top\mathrm{D}^2\eta(U)R_j(U)>0,\qquad j=1,\cdots,n.
\label{align:ch07-熵凸}
\end{align}

当方程组\eqref{align:ch07-典范守恒律}是对称的，
\begin{align}
\mathrm{D}F(U)^\top=\mathrm{D}F(U),\qquad U\in\mathscr{O},
\label{align:ch07-对称双曲}
\end{align}
它容许两个有趣的熵--熵流对：
\begin{align}
\eta=\tfrac12|U|^2,\qquad q=U\cdot F(U)-h(U),
\label{align:ch07-熵对1}
\end{align}
\begin{align}
\eta=h(U),\qquad q=\tfrac12|F(U)|^2,
\label{align:ch07-熵对2}
\end{align}
其中 $h$ 由条件
\begin{align}
\mathrm{D}h(U)=F(U)^\top
\label{align:ch07-h势}
\end{align}
定义。

如第 III 章所解释，方程组\eqref{align:ch07-热弹性}、\eqref{align:ch07-p系统}、\eqref{align:ch07-理想气体}分别配有熵--熵流对
\begin{align}
\eta=-s,\qquad q=0,
\label{align:ch07-热弹性熵}
\end{align}
\begin{align}
\eta=\tfrac12 v^2+e(u),\qquad q=-v\sigma(u),\qquad e(u)=\int^u\sigma(\omega)\,\mathrm{d}\omega,
\label{align:ch07-p系统熵}
\end{align}
\begin{align}
\eta=\frac{1}{2}\rho v^2+\frac{\kappa}{\gamma-1}\rho^\gamma,\qquad
q=\frac{1}{2}\rho v^3+\frac{\kappa\gamma}{\gamma-1}\rho^\gamma v,
\label{align:ch07-理想气体熵}
\end{align}
它们由热力学第二定律诱导。\footnote{在等温情形 $\gamma=1$，方程组\eqref{align:ch07-理想气体}的熵--熵流对取如下形式：$\eta=\frac12\rho v^2+\kappa\rho\log\rho$，$q=\frac12\rho v^3+\kappa\rho v\log\rho+\kappa\rho v$。}事实上，把 $v\sigma$ 与 $\sigma\mathrm{d}\omega$ 解释为 $v\cdot\sigma$ 与 $\sigma\cdot\mathrm{d}\omega$，\eqref{align:ch07-p系统熵}即使对方程组\eqref{align:ch07-平面振荡}也构成一个熵--熵流对。当表示为典范状态变量（即\eqref{align:ch07-热弹性}的 $(u,v,E)$、\eqref{align:ch07-p系统熵}的 $(u,v)$、\eqref{align:ch07-理想气体熵}的 $(\rho,m)$）的函数时，上述熵是凸的。

在发展方程组\eqref{align:ch07-典范守恒律}的理论时，构造具有给定规格的熵将是有用的。这些熵必须是\eqref{align:ch07-可积性}的解，而\eqref{align:ch07-可积性}是关于单个未知量 $\eta$ 的、含 $\frac12 n(n-1)$ 个偏微分方程的线性二阶方程组。于是当 $n=2$ 时，\eqref{align:ch07-可积性}退化为单个线性双曲方程，可求解以产生大量熵。相反，对 $n\geqslant 3$，\eqref{align:ch07-可积性}形式上是过定的。尽管存在如\eqref{align:ch07-熵对1}与\eqref{align:ch07-熵对2}这样的特殊解，除非\eqref{align:ch07-典范守恒律}是特殊的，否则不应指望熵大量存在。值得注意的是，当\eqref{align:ch07-典范守恒律}配有 Riemann 不变量坐标系 $(w_1,\cdots,w_n)$ 时，\eqref{align:ch07-可积性}的过定性消失。在此情形下，把 $\eta$ 与 $q$ 作为状态向量 $W=(w_1,\cdots,w_n)^\top$ 的函数来寻求是方便的。从右乘 $R_j(U)$ 于\eqref{align:ch07-熵流对定义}并利用\eqref{align:ch07-链式法则}，我们推出\eqref{align:ch07-熵流对定义}现在等价于
\begin{align}
\frac{\partial q}{\partial w_j}=\lambda_j\frac{\partial\eta}{\partial w_j},\qquad j=1,\cdots,n.
\label{align:ch07-熵流W}
\end{align}
与\eqref{align:ch07-熵流W}相关联的可积性条件取形式
\begin{align}
\frac{\partial^2\eta}{\partial w_j\partial w_k}+g_{jk}\frac{\partial\eta}{\partial w_j}+g_{kj}\frac{\partial\eta}{\partial w_k}=0,\qquad j,k=1,\cdots,n;\ j\neq k,
\label{align:ch07-熵可积W}
\end{align}
其中 $g_{jk},g_{kj}$ 是通过\eqref{align:ch07-gjk定义}定义的函数。如果由\eqref{align:ch07-熵可积正交}直接推导\eqref{align:ch07-熵可积W}，就出现 $g_{jk}$ 的另一个有用表达式。事实上，对 $j,k=1,\cdots,n$，
\begin{align}
R_j^\top\mathrm{D}^2\eta R_k=\mathrm{D}(\mathrm{D}\eta R_j)R_k-\mathrm{D}\eta\mathrm{D}R_j R_k
=\mathrm{D}(\mathrm{D}\eta R_j)R_k-\sum_{i=1}^n\frac{\partial\eta}{\partial w_i}\mathrm{D}w_i\mathrm{D}R_j R_k.
\label{align:ch07-D2eta展开}
\end{align}
合并\eqref{align:ch07-熵可积正交}、\eqref{align:ch07-链式法则}、\eqref{align:ch07-Lie括号零}与\eqref{align:ch07-DwDR}，我们得到形如\eqref{align:ch07-熵可积W}的方程，其中
\begin{align}
g_{jk}=R_j^\top\mathrm{D}^2 w_j R_k,\qquad j,k=1,\cdots,n;\ j\neq k.
\label{align:ch07-gjk二阶}
\end{align}

读者可借助\eqref{align:ch07-Lie括号恒等式}、\eqref{align:ch07-Riemann规范化}、\eqref{align:ch07-D2w交叉}、\eqref{align:ch07-Lie括号零}、\eqref{align:ch07-DwDR}与\eqref{align:ch07-链式法则}作为习题直接验证：\eqref{align:ch07-gjk定义}与\eqref{align:ch07-gjk二阶}等价。

取 $k=j$，利用\eqref{align:ch07-链式法则}、\eqref{align:ch07-DwDR}并回忆\eqref{align:ch07-熵凸}，应用\eqref{align:ch07-D2eta展开}，我们推出：以 Riemann 不变量表示，$\eta$ 的凸性条件由不等式组
\begin{align}
\frac{\partial^2\eta}{\partial w_j^2}+\sum_{i=1}^n a_{ij}\frac{\partial\eta}{\partial w_i}\geqslant 0,\qquad j=1,\cdots,n,
\label{align:ch07-熵凸W}
\end{align}
表达，其中
\begin{align}
a_{ij}=R_j^\top\mathrm{D}^2 w_i R_j,\qquad i,j=1,\cdots,n.
\label{align:ch07-aij}
\end{align}

方程组\eqref{align:ch07-熵可积W}关于单个未知量 $\eta$ 含 $\frac12 n(n-1)$ 个方程，因而当 $n\geqslant 3$ 时似乎是过定的。然而结果表明，该方程组内部是相容的。为看出这一点，对 $w_i$ 求\eqref{align:ch07-熵可积W}的导数，得到
\begin{align}
\frac{\partial^3\eta}{\partial w_i\partial w_j\partial w_k}
=-\frac{\partial g_{jk}}{\partial w_i}\frac{\partial\eta}{\partial w_j}
+g_{jk}\left(g_{ji}\frac{\partial\eta}{\partial w_j}+g_{ij}\frac{\partial\eta}{\partial w_i}\right)
-\frac{\partial g_{kj}}{\partial w_i}\frac{\partial\eta}{\partial w_k}
+g_{kj}\left(g_{ki}\frac{\partial\eta}{\partial w_k}+g_{ik}\frac{\partial\eta}{\partial w_i}\right).
\label{align:ch07-熵三阶导}
\end{align}

方程组\eqref{align:ch07-熵可积W}可积，当且仅当对 $i\neq j\neq k\neq i$，\eqref{align:ch07-熵三阶导}的右端关于 $(i,j,k)$ 对称。但由恒等式\eqref{align:ch07-g对称}与\eqref{align:ch07-gCocycle}，这总是成立。因而，可以在任一给定状态 $\bar{W}=(\bar{w}_1,\cdots,\bar{w}_n)^\top$ 的邻域内构造熵 $\eta$，使其沿平行于坐标轴的直线上取预值 $\{\eta(w_1,\bar{w}_2,\cdots,\bar{w}_n),\eta(\bar{w}_1,w_2,\cdots,\bar{w}_n),\cdots,\eta(\bar{w}_1,\cdots,\bar{w}_{n-1},w_n)\}$。当 $n=2$ 时，这相当于求解一个经典 Goursat 问题。

我们由此证明：配有 Riemann 不变量坐标系的方程组也配有丰富的熵。正因如此，这类方程组被称为\emph{丰富的（rich）}。特别地，电泳方程组\eqref{align:ch07-电泳}与电磁波方程组\eqref{align:ch07-电磁波}是丰富的。读者可在\cref{section:ch07-注记}
所引文献中找到如何构造这些方程组的熵族。

\end{document}
